\documentclass[conference]{../../../../Conference-LaTeX-template_10-17-19/IEEEtran}
\usepackage{amsmath,amssymb,amsfonts,array,booktabs,tabularx,balance,url}
\newcolumntype{Y}{>{\raggedright\arraybackslash}X}
\newcommand{\R}{\mathbb{R}}

\title{Will the Three-Dimensional Navier--Stokes Problem Be Solved?\\
An Audit-First Research Program for Critical-Scale Regularity}
\author{\IEEEauthorblockN{Research Proposal}
\IEEEauthorblockA{Evidence-bounded four-agent synthesis -- Mathematical design only}}
\begin{document}
\maketitle

\begin{abstract}
The three-dimensional incompressible Navier--Stokes existence and smoothness problem asks whether every smooth, divergence-free, finite-energy initial velocity generates a unique smooth solution for all positive time, or whether a finite-time singularity can occur. This is a continuum-PDE question, not a request for another successful numerical simulation. The surveyed record contains Leray weak solutions, partial regularity, local smooth theory, critical continuation criteria, quantitative concentration results, and special-class global regularity theorems. It does not contain a cross-validated, conventionally accepted complete resolution. Recent public manuscripts declaring a solution are therefore classified as claims under audit rather than as closure evidence.

The selected direction, ScaleBridge, targets the exact point at which known estimates fail: finite energy does not unconditionally control a scale-critical norm. It formulates a pressure-safe critical-envelope program around a hypothetical first singularity. The program uses local energy and pressure identities, epsilon-regularity, concentration extraction, vortex-stretching analysis, compactness, and ancient-solution rigidity. Its decisive open obligation is a uniform non-circular decay estimate for a scale-invariant velocity-pressure-stretching envelope. If that estimate is proved and its nonlocal tail is absorbed without importing global regularity, iteration reaches the epsilon threshold and excludes the singularity. If it fails, the design produces a typed obstruction rather than an overstated proof claim. Mathematical derivations are research obligations only; no computation, numerical experiment, proof-assistant run, singularity, or new theorem is claimed.
\end{abstract}
\begin{IEEEkeywords}
Navier--Stokes equations, global regularity, Millennium Prize problem, critical spaces, vortex stretching, partial regularity, proof audit, mathematical fluid dynamics
\end{IEEEkeywords}

\section{Introduction}
The incompressible Navier--Stokes equations are central to mathematical fluid dynamics because they combine a dissipative linear term with a nonlinear transport term strong enough to create multiscale turbulent behavior. They are used throughout science and engineering, but their practical usefulness does not decide the associated theorem. The unresolved three-dimensional question is whether a classical solution can develop a finite-time singularity from smooth data, or whether viscosity always prevents that event. A proof must settle an all-data, all-time, all-scale alternative, not merely show that a calculation on a finite mesh remains stable.

For velocity $u:\R^3\times[0,T)\to\R^3$, pressure $p$, and viscosity $\nu>0$, the Cauchy problem is
\begin{equation}
\partial_t u+(u\cdot\nabla)u+\nabla p=\nu\Delta u,\qquad
\nabla\cdot u=0,\qquad u(\cdot,0)=u_0.
\label{eq:NS}
\end{equation}
The Millennium formulation gives whole-space and periodic versions under carefully stated assumptions \cite{fefferman}. A positive resolution proves global smooth existence and uniqueness for every admissible smooth divergence-free datum. A negative resolution constructs an admissible datum for which a singularity forms in finite time. It is not enough to alter the PDE with a molecular cutoff, restrict the datum to an axisymmetric class, or replace the continuum solution by a discretization.

The report answers two questions. First, what is justified by the surveyed evidence? The direct answer is that the general $3$D problem remains open in the evidence record used here. Second, what kind of research program could actually advance it? The answer is ScaleBridge, an audit-first architecture that transforms a prospective proof into explicit analytic obligations. Its aim is to expose the missing bridge between energy-level control and scale-critical regularity, and thereby make both valid progress and invalid claims recognizable.

The four-agent workflow is used in a disciplined sense. The Survey stage freezes the theorem record and the research gaps. The Idea stage compares routes and selects a pressure-safe concentration program. The ExperimentDesign stage treats a proof attempt as a mathematical design with variables, decision rules, and counterexample protocols. The Author stage writes only within those inputs. The report has proposal status and does not claim that its displayed derivations solve \eqref{eq:NS}.

\section{Mathematical Statement and Established Structure}
\subsection{Scaling and the energy barrier}
Equation \eqref{eq:NS} is invariant under the parabolic scaling
\begin{equation}
u_\lambda(x,t)=\lambda u(\lambda x,\lambda^2t),\qquad
p_\lambda(x,t)=\lambda^2p(\lambda x,\lambda^2t).
\label{eq:scaling}
\end{equation}
The scaling identifies borderline quantities. The spatial $L^3$ norm of $u$ and the spacetime norms satisfying
\begin{equation}
\frac{2}{q}+\frac{3}{r}=1
\label{eq:critical}
\end{equation}
are invariant when $u\in L_t^qL_x^r$. A norm that becomes smaller under rescaling is subcritical; a norm that grows is supercritical. This classification determines whether an a priori estimate plausibly controls the fine scales where a singularity would occur.

For smooth rapidly decaying solutions, multiplying \eqref{eq:NS} by $u$ and using incompressibility gives
\begin{equation}
\frac{1}{2}\frac{d}{dt}\|u(t)\|_{L^2}^2+
\nu\|\nabla u(t)\|_{L^2}^2=0.
\label{eq:energy}
\end{equation}
Leray used the corresponding inequality to construct global finite-energy weak solutions \cite{leray}. Equation \eqref{eq:energy} controls total kinetic energy and viscous dissipation, but the $L^2$ norm is not critical for the three-dimensional scaling. A hypothetical concentration can preserve the energy budget while becoming increasingly intense at smaller radii. Therefore no algebraic rewriting of \eqref{eq:energy} alone proves global smoothness.

Taking curl gives the vorticity equation
\begin{equation}
\partial_t\omega+(u\cdot\nabla)\omega-\nu\Delta\omega
 =(\omega\cdot\nabla)u,\qquad \omega=\nabla\times u.
\label{eq:vorticity}
\end{equation}
The right side is vortex stretching. Formally pairing \eqref{eq:vorticity} with $\omega$ gives
\begin{equation}
\frac{1}{2}\frac{d}{dt}\|\omega(t)\|_{L^2}^2+
\nu\|\nabla\omega(t)\|_{L^2}^2
=\int_{\R^3}(\omega\cdot\nabla u)\cdot\omega\,dx.
\label{eq:enstrophy}
\end{equation}
In two dimensions the analogous stretching term vanishes. In three dimensions the term in \eqref{eq:enstrophy} has no general sign. A real solution must establish an inequality that controls it at every relevant scale, rather than describe a favorable average tendency.

\subsection{What continuation criteria prove}
Known regularity criteria provide conditional theorems. A Prodi--Serrin statement says that a smooth solution can be continued if a critical or subcritical spacetime norm satisfying \eqref{eq:critical} is finite. The endpoint result of Escauriaza, Seregin, and Sverak shows that boundedness in the critical space $L_t^\infty L_x^3$ prevents a finite blow-up in its stated setting \cite{ess}. These theorems are exact implications:
\begin{equation}
\sup_{0<t<T}\|u(t)\|_{L^3}<\infty
\quad\Longrightarrow\quad
\text{continuation past }T.
\label{eq:continuation}
\end{equation}
The unresolved issue is the antecedent of \eqref{eq:continuation}. A proof of global regularity must derive such a bound, or a comparable continuation condition, for arbitrary smooth data without assuming the desired conclusion. Applying \eqref{eq:continuation} after postulating the critical bound simply moves the open problem. Critical-space analyses using concentration compactness and rigidity clarify this point \cite{kenigkoch,tao}.

There are global theorems for small data, special symmetries, thin domains, strong rotation, and well-prepared large data. These results reveal mechanisms that can suppress nonlinear growth, but their extra assumptions matter. A theorem for a structured subclass cannot be extended by rhetoric to all smooth three-dimensional flows. The exact domain, boundary condition, initial-data class, and norm are part of every valid theorem statement.

\subsection{Partial regularity and the singularity alternative}
Suitable weak solutions obey a localized energy inequality. Caffarelli, Kohn, and Nirenberg proved partial regularity, showing that possible singular points form a small parabolic set \cite{ckn}. The result does not say that the set is empty. Its importance for a future proof is that it yields epsilon-regularity principles: if a normalized local velocity-pressure quantity is small in a cylinder, the solution is regular in a smaller cylinder.

Let $z_0=(x_0,t_0)$ and
\begin{equation}
Q_r(z_0)=B_r(x_0)\times(t_0-r^2,t_0).
\end{equation}
Two standard scale-invariant quantities are
\begin{align}
C(r,z_0)&=r^{-2}\int_{Q_r(z_0)}|u|^3\,dx\,dt, \label{eq:C}\\
D(r,z_0)&=r^{-2}\int_{Q_r(z_0)}
|p-(p)_{B_r}(t)|^{3/2}\,dx\,dt. \label{eq:D}
\end{align}
The exact epsilon-regularity formulation depends on the solution class and normalization. Its logical form is
\begin{equation}
C(r,z_0)+D(r,z_0)<\varepsilon_\star
\quad\Longrightarrow\quad
z_0\text{ is regular}.
\label{eq:epsilon}
\end{equation}
Thus a first singular point would have to sustain non-small critical concentration along a sequence of radii. The global problem may be reframed as follows: either prove that this sustained concentration is impossible, or construct it in an admissible solution.

\section{Survey Status and Claim-Audit Boundary}
\subsection{Evidence-supported status}
The survey registry separates theorem inputs from motivation. Leray weak existence, partial regularity, endpoint continuation, ancient-solution Liouville results, and localization equivalences are PROVED inputs within their stated hypotheses \cite{leray,ckn,ess,knss,tao}. Quantitative work has sharpened the constraints a hypothetical blow-up must satisfy. Barker and Prange obtain quantitative spatial-concentration information under declared critical-norm assumptions \cite{barkerprange}. Farwig gives a modern survey that places these advances in the historical development from Leray to the Millennium problem \cite{farwig}.

The status decision is not a prediction about whether a solution will appear soon. Mathematics provides no justified timetable. A claim that the problem will be solved becomes meaningful only when it is translated into a theorem route: an all-data a priori estimate, a blow-up construction, a compactness-and-rigidity contradiction, or another exact mechanism. The evidence supports vigorous research, but it does not license a forecast that a final theorem is imminent.

Recent numerical work reports apparent self-attenuation of extreme turbulent events through a local/nonlocal strain decomposition \cite{buaria}. This is scientifically interesting because it suggests a depletion mechanism. It is not evidence of a global PDE theorem. Simulations have finite resolution; the conclusion must survive every scale, every smooth datum, and every pressure configuration. The present proposal uses the observation only to motivate a testable analytic inequality.

\subsection{How public solution declarations are treated}
The dual scholarly search located public 2025--2026 records that describe themselves as resolutions of the problem. A literature search is a discovery procedure, not a proof validator. The records are labeled \texttt{CLAIM\_UNDER\_AUDIT}. This label does not say that a manuscript is false. It says that its theorem chain is not an admissible premise in this proposal.

Every claimed solution must pass the same audit:
\begin{enumerate}
\item Does it state the exact unrestricted equation, domain, data class, and solution class required by the target problem?
\item Does every new quantity have a defined norm, scaling law, and domain of validity?
\item Is a claimed global bound derived from the equation, rather than imposed by a new physical axiom, finite discretization, or altered admissibility class?
\item Are pressure reconstruction, compactness, limits, and uniqueness handled in the same class used by the conclusion?
\item Does the final implication cover all data and all times, rather than an initially assumed bounded envelope?
\end{enumerate}
A missing answer leaves an open obligation. A finite-grid cutoff cannot settle a continuum theorem. A new axiom that rules out arbitrarily small scales may describe a different physical model, but it cannot be inserted into the classical Navier--Stokes formulation without proving that it follows from \eqref{eq:NS}. The user-selected in-app browser failed before it could reach official Clay or publisher pages, so direct page verification is a human-review requirement.

\begin{table}[t]
\caption{Status Labels Used in the ScaleBridge Audit}
\centering\small
\begin{tabular}{p{0.27\columnwidth}p{0.54\columnwidth}}
\toprule
Label & Meaning\\
\midrule
PROVED & A theorem with stated hypotheses and a verifiable source.\\
CONDITIONAL & A valid implication with an explicitly retained assumption.\\
FINITE & A calculation or estimate valid only for a declared range.\\
HEURISTIC & A numerical or conceptual guide without theorem force.\\
OPEN OBLIGATION & A missing estimate, map, or limiting argument.\\
CLAIM UNDER AUDIT & A public declaration not yet accepted as a proof input.\\
\bottomrule
\end{tabular}
\end{table}

\section{Why Energy Does Not Close the Problem}
\subsection{The critical-scale gap}
The energy identity \eqref{eq:energy} controls
\begin{equation}
u\in L_t^\infty L_x^2\cap L_t^2\dot H_x^1,
\label{eq:energyclass}
\end{equation}
while the continuation criterion \eqref{eq:continuation} asks for a scale-critical quantity. Interpolation produces useful spacetime integrability, but not a known unconditional endpoint estimate strong enough to prohibit every concentration cascade. The word supercritical records this mismatch between the available estimate and the scaling of a prospective singularity.

Suppose a smooth solution exists on $[0,T_\star)$ but cannot be continued at $T_\star$. Any valid regularity criterion must fail as $t\uparrow T_\star$. This does not mean all norms diverge at the same rate or location. It means a proof cannot use boundedness of its favorite critical norm unless it derives that boundedness independently. Quantitative blow-up results refine necessary behavior under additional norm hypotheses \cite{barkerprange}; they do not establish those hypotheses for arbitrary flows.

Rescaling around a candidate concentration point is attractive. It can turn a sequence of small cylinders into a normalized solution on a fixed cylinder or an ancient time interval. But rescaling creates three linked requirements: compactness of the sequence, nontriviality of the limit, and rigidity of the exact limit class. A trivial limit says nothing; a Liouville theorem outside the limit class says nothing; a compactness theorem requiring the bound under investigation is circular.

\subsection{Pressure and nonlocal strain}
Incompressibility ties pressure to velocity through
\begin{equation}
-\Delta p=\partial_i\partial_j(u_i u_j).
\label{eq:pressure}
\end{equation}
Consequently, a local cylinder analysis has a nonlocal component. A cutoff produces flux and commutator terms. A pressure decomposition must separate local and harmonic or far-field pieces, then estimate each in a compatible scale-invariant norm. A proposed local inequality that deletes the nonlocal term is not a localization proof.

Let
\begin{equation}
S=\frac{1}{2}(\nabla u+\nabla u^\mathsf{T})
\end{equation}
be the strain tensor. Then
\begin{equation}
(\omega\cdot\nabla u)\cdot\omega=\omega\cdot S\omega.
\label{eq:stretch}
\end{equation}
At a candidate singular scale, define
\begin{equation}
\mathcal V(r,z_0)=r\int_{Q_r(z_0)}
(\omega\cdot S\omega)_+\,dx\,dt.
\label{eq:V}
\end{equation}
The factor $r$ makes \eqref{eq:V} scale invariant. It records the part of stretching that can amplify enstrophy; it does not estimate pressure, the negative part, or dissipation. Any coercivity result must connect \eqref{eq:V} to localized dissipation while retaining the far-field contribution.

\subsection{What a genuine bridge would look like}
A genuine advance could give an all-data critical estimate, rule out a complete class of normalized ancient profiles, or construct a valid singularity. It must cross the relevant quantifier gap:
\begin{equation}
\text{finite energy and smooth data}
\Longrightarrow
\text{uniform critical envelope}
\Longrightarrow
\text{continuation for all }T.
\label{eq:bridge}
\end{equation}
The first implication in \eqref{eq:bridge} is not present in the established record. ScaleBridge narrows it to a local version near a hypothetical first singularity. That replacement is legitimate only if local decay iterates over all scales and covers every possible concentration point.

\section{ScaleBridge: Selected Research Direction}
\subsection{Research objective}
ScaleBridge asks whether a pressure-safe, scale-invariant estimate can force the critical concentration associated with a first hypothetical singularity to decay below an epsilon-regularity threshold. This objective names state variables, scale, decision threshold, and outcomes. A successful theorem excludes the first singularity. A failed inequality identifies whether the obstruction is stretching, pressure, far-field strain, compactness, or an incompatible rigidity class.

Define a proposal envelope
\begin{equation}
\mathcal E(r,z_0)=
a\,C(r,z_0)+b\,D(r,z_0)+c\,\mathcal V(r,z_0)+\mathcal T(r,z_0),
\label{eq:envelope}
\end{equation}
where $a,b,c>0$ are fixed only after a derivation, and $\mathcal T$ records cutoff, exterior-strain, and pressure-tail terms. The eventual proof may not choose weights after observing a favorable example. It must state weights, radius range, solution class, and constants before iteration begins.

The target is
\begin{equation}
\mathcal E(\theta r,z_0)\le q\mathcal E(r,z_0)
+K\mathcal T_{\mathrm{rem}}(r,z_0),
\qquad 0<\theta,q<1,
\label{eq:decay}
\end{equation}
where the remainder is summable or absorbed by a proved estimate that does not assume global regularity. If \eqref{eq:decay} is uniform at all candidate cylinders, iteration drives $\mathcal E$ below $\varepsilon_\star$. Compatibility with \eqref{eq:epsilon} then rules out the candidate singularity.

\subsection{Why this route is discriminating}
The route has a theorem endpoint and an obstruction endpoint. The theorem endpoint requires every edge in the proof graph to be PROVED. The obstruction endpoint may show that an alignment hypothesis fails to control exterior strain, or that a pressure estimate is scale-supercritical. Such an obstruction does not disprove Navier--Stokes regularity. It prevents a particular proof program from being misreported as a resolution.

The choice protects against a recurrent error in public proof claims: replacing a continuum all-scale question with finite information. A calculation can evaluate $C(r,z_0)$ down to a finite radius, but it cannot show \eqref{eq:decay} at all radii. A simulation can reveal local alignment, but it cannot establish a universal $q<1$. A proof based on a fixed microscopic cutoff solves a regularized model unless it first proves that the original PDE supplies that cutoff.

\subsection{Portfolio alternatives}
The idea stage compared ScaleBridge with two neighboring programs. The first, critical-element compactness, seeks a minimal blow-up profile and a Liouville contradiction. It is competitive and partly integrated into O7, but it needs a rigidity theorem matching the extracted profile. The second, multiscale alignment, seeks direct geometric cancellation in \eqref{eq:stretch}. It is high-risk because it may discover the decisive estimate, but it must control local and nonlocal strain in every orientation. A finite-grid proof and unreviewed declaration route were rejected because they violate the problem definition or evidence policy.

\section{Experiment Design: Formal Derivation and Gates}
\subsection{Design-only mathematical experiment}
Here an experiment means a future formal derivation, certified estimate, or adversarial test of a lemma. No execution occurs in this report. There are no observed numerical results, no DNS output, no proof-assistant certificates, and no newly computed solutions. The design specifies what a future result must contain before it may change the problem status.

The first derivation begins with a smooth solution and a nonnegative cutoff $\phi$ supported in $Q_r(z_0)$. Multiplication of \eqref{eq:NS} by $u\phi^2$ produces a local energy balance containing viscosity, transport flux, pressure flux, and cutoff derivatives. Taking curl and localizing \eqref{eq:vorticity} produces an enstrophy balance containing \eqref{eq:stretch}. These identities define the terms that the envelope \eqref{eq:envelope} must pay for.

The second derivation reconstructs pressure through \eqref{eq:pressure} and splits it into a local singular-integral term and a harmonic or far-field term. The scale of each is checked against \eqref{eq:C}--\eqref{eq:D}. A candidate proof is rejected if it estimates a harmonic remainder by a global quantity unavailable at the first singularity.

\begin{table*}[t]
\caption{ScaleBridge Proof Obligations and Release Gates}
\centering\small
\begin{tabularx}{\textwidth}{p{0.07\textwidth}Y Y}
\toprule
ID & Required result & Release gate and failure mode\\
\midrule
O1 & Exact formulation, local theory, scaling, and continuation statement. & Domain, data class, pressure normalization, and solution class must match the conclusion.\\
O2 & Localized energy, vorticity, and pressure identities with cutoff and far-field terms. & No flux, pressure, commutator, or exterior strain term may be discarded.\\
O3 & Epsilon-regularity in the normalization of $C$, $D$, and $\mathcal E$. & A differently normalized theorem needs a proved conversion.\\
O4 & First-concentration extraction from a hypothetical singularity. & The selected sequence remains divergence free, scale normalized, and nontrivial.\\
O5 & Uniform pressure-safe vortex-stretching depletion giving \eqref{eq:decay}. & A numerical trend, fixed radius, or conditional bound is insufficient.\\
O6 & Tail absorption or summability for $\mathcal T_{\mathrm{rem}}$. & The argument may not assume the critical norm it seeks.\\
O7 & Iteration to epsilon regularity, or a compact ancient profile in the exact rigidity class. & A trivial limit or unmatched Liouville theorem blocks contradiction.\\
O8 & Composition and independent proof audit. & Every prior edge must be PROVED; a GAP blocks a global claim.\\
\bottomrule
\end{tabularx}
\end{table*}

\subsection{Derivation of the concentration alternative}
Assume for contradiction that $T_\star$ is a first singular time. The local implication \eqref{eq:epsilon} gives a contraposition: every sufficiently small neighborhood of a singular point has non-small critical concentration. Thus one can select radii $r_k\downarrow0$ and centers $z_k$ with a non-small normalized envelope. The design does not claim that selection provides a limiting profile. It creates the input for O4.

Rescale around $(z_k,r_k)$:
\begin{equation}
u^{(k)}(y,s)=r_k u(x_k+r_ky,t_k+r_k^2s),\qquad
p^{(k)}(y,s)=r_k^2p(x_k+r_ky,t_k+r_k^2s).
\label{eq:rescale}
\end{equation}
Equation \eqref{eq:NS} and the quantities in \eqref{eq:C}--\eqref{eq:V} have compatible scalings. A compactness argument must name the topology in which $u^{(k)}$ and $p^{(k)}$ converge, prove sufficient uniform bounds, preserve the local energy inequality, and prove that the limit is nonzero. Strong convergence is not granted merely because the sequence is visually similar after rescaling.

If the decay lemma \eqref{eq:decay} is available, a compactness route might become unnecessary: repeated shrinkage yields a geometric series plus a controlled remainder. If it is unavailable, a critical-element route attempts to pass from \eqref{eq:rescale} to a bounded ancient solution. Liouville theorems for selected ancient classes are powerful inputs \cite{knss}, but their hypotheses must be checked line by line. The design prohibits the assertion that every ancient solution vanishes, because available theorems have particular boundedness, symmetry, or admissibility conditions.

\subsection{The decisive new lemma}
O5 seeks a local inequality with a strict gain:
\begin{equation}
\mathcal V(\theta r,z_0)
\le \eta\,\mathcal D(r,z_0)
+A\bigl(C(r,z_0)+D(r,z_0)\bigr)
+B\mathcal T_{\mathrm{far}}(r,z_0),
\qquad \eta<1,
\label{eq:depletion}
\end{equation}
where $\mathcal D$ is localized viscous dissipation in the chosen normalization. The coefficients must be independent of the first-singularity scale. An estimate of the form \eqref{eq:depletion} would convert stretching into a perturbative part of the local energy inequality. Combined with a pressure estimate and interpolation, it could yield \eqref{eq:decay}.

Equation \eqref{eq:depletion} is a target, not a result. It may fail because local strain opposes amplification while the far field supplies an uncontrolled positive component, because alignment holds in a model class but not for arbitrary smooth solutions, because $\eta$ approaches one along an adversarial sequence, or because the tail is scale critical and non-summable. Each failure identifies needed structure and prevents a heuristic from being promoted to a theorem.

\section{Validation, Counterexamples, and Claim Audit}
\subsection{Mathematical validation protocol}
The validation protocol is stronger than checking whether algebraic lines look plausible. Every lemma has a typed input and output. For O2, the input is a smooth or suitable weak solution plus a declared cutoff; the output is a local identity with all terms recorded. For O5, the input is the exact local class and normalized quantities; the output is an inequality with constants, radius range, and dependence stated. For O7, the input is the compactness sequence; the output is epsilon regularity or an ancient profile satisfying the hypotheses of a cited rigidity theorem.

An all-data theorem cannot be assembled from statements with incompatible quantifiers. For example,
\begin{equation}
\forall r\ge r_0,\quad \mathcal E(\theta r)\le q\mathcal E(r)+R(r)
\label{eq:allradii}
\end{equation}
is different from verifying the same inequality at finitely many radii. Likewise, a result for a special symmetry class cannot be invoked after a general rescaling unless the symmetry is inherited. The review record for each edge lists domain, norm, time interval, and whether it is unconditional.

Proof-assistant formalization could later strengthen the audit of O1--O4 and the iteration in O7. It would not produce O5. A formal environment checks that a stated derivation follows from imported theorems; it cannot decide whether the imported analytic lemma is true unless that lemma is formally established. The proposal treats formalization as an anti-ambiguity tool, not an automatic solution machine.

\subsection{Counterexample and obstruction protocol}
The program recognizes two counterexample notions. A counterexample to ScaleBridge is a normalized admissible sequence for which \eqref{eq:depletion} cannot hold. It is an obstruction to this strategy, not a singular solution of Navier--Stokes. A counterexample to the Millennium problem requires a smooth initial datum whose solution develops a finite-time singularity in the specified formulation. Confusing the two would turn a failed technique into an invalid negative resolution.

For a putative ScaleBridge obstruction, the audit asks: Is the object a solution of the original PDE or a numerical surrogate? Does it have the stated pressure normalization? Is the scale $r$ allowed? Are the quantities in \eqref{eq:C}, \eqref{eq:D}, and \eqref{eq:V} finite and correctly normalized? Is the far-field term included? Does the proposed sequence obey every hypothesis? Only then can the result be tagged as an obstruction.

For a claimed positive solution, the audit draws a dependency graph from smooth data to a global solution. Each arrow records an exact theorem, all hypotheses, constants, and quantifiers. Any arrow that says "the small scales are bounded by physics," "the numerical grid resolves all scales," or "the envelope stays finite by construction" is a GAP unless it is derived from \eqref{eq:NS}. The procedure judges whether the displayed implication reaches the mathematical target.

\subsection{Interpretation of mixed outcomes}
Mixed outcomes are useful. A proof of O2 without O5 is a correct local identity, not global regularity. A proof of O5 under an alignment assumption is a conditional theorem, not the all-data result. A numerical experiment suggesting that $\mathcal V$ is small at a selected Reynolds number is a heuristic. A compactness construction with a profile outside a known Liouville class is a research lead. The release label changes with the evidence; valid subordinate results are retained.

This approach is not overconservative. A uniform decay lemma with verified tail absorption would be a major advance, and a complete O1--O8 chain would be a candidate solution deserving immediate expert review. The discipline lies only in refusing to name a result before its hypotheses and universal quantifiers have been discharged.

\section{Risks, Limitations, and Human Review}
The central technical risk is that the desired depletion inequality \eqref{eq:depletion} may be false for arbitrary smooth flows. Vorticity geometry can be favorable in many regimes while failing in an adversarial cascade. A second risk is that a local pressure estimate creates a nonlocal tail at critical scaling, preventing iteration. A third is that an ancient-profile limit may not lie in any available rigidity class. None of these risks establishes a singularity; each marks a place where the selected route can fail.

The source limitations are narrower. The project conducted dual-engine discovery through OpenAlex and AnySearch and froze metadata for the cited record. The explicit in-app browser requested for this work was blocked by a Windows access-control failure before official webpages could be inspected. A mathematical-fluid-dynamics reviewer should refresh official Clay status, DOI metadata, versions, and recognition records from publisher or institute pages before public dissemination. This is a source-verification limitation, not evidence for a solution or against one.

The report requires independent expert review before a future theorem claim. One reviewer should inspect local PDE estimates and pressure decomposition; another should examine the compactness or rigidity chain; and, if computation is used, an independent implementation should verify finite certificates. No amount of source fingerprinting replaces a line-by-line proof audit. The browser limitation does not prevent the present bounded conclusion: no cross-validated accepted complete proof was found in the survey record.

\section{Conclusion}
Will the Navier--Stokes problem ever be solved? The evidence cannot supply a calendar prediction. It supplies a sharper statement: the problem will be solved only when an argument proves global smoothness for every datum in the Millennium class or constructs a valid singularity in that class. Current evidence supports major partial results--weak existence, partial regularity, critical continuation criteria, quantitative concentration bounds, and special-case global theorems--but it does not support declaring the general $3$D problem resolved.

ScaleBridge contributes an architecture for the missing step. It localizes a first hypothetical singularity, defines a scale-invariant velocity-pressure-stretching envelope, and identifies a pressure-safe strict decay estimate as the decisive open lemma. If that lemma and tail absorption are proved non-circularly, epsilon-regularity excludes the singularity. If a compact profile survives instead, the program requires an exact rigidity theorem before a contradiction is announced. If either route fails, the output is an obstruction ledger rather than an inflated claim.

The result is a direct research proposal: not a declaration that global regularity has been obtained, but a rigorous way to decide which future estimates would move the problem toward a solution. That distinction preserves the force of known theorems while concentrating future work on the universal, critical-scale bridge they do not yet provide.

\appendices
\section{Formal Proof-Obligation Graph}
The ScaleBridge graph begins at the exact formulation of \eqref{eq:NS}. Its first nodes are smooth local existence, scaling, the energy identity, the local energy inequality, and epsilon regularity. These are not interchangeable: a statement about smooth solutions is not automatically a statement about suitable weak limits, and a whole-space pressure normalization may not transfer to a periodic proof without a separate argument.

The graph then branches at a hypothetical first singularity. One branch selects concentration cylinders, derives \eqref{eq:decay}, absorbs the tail, and iterates to \eqref{eq:epsilon}. The other branch rescales by \eqref{eq:rescale}, establishes compactness and nontriviality, and invokes a rigidity theorem matched to the resulting ancient solution. Both branches terminate at continuation, but neither may borrow an edge from the other unless its hypotheses have been proved.

Every node has a label from Table I. A PROVED node may support a later edge. A CONDITIONAL node carries its condition forward. A FINITE node may support only a finite conclusion. A HEURISTIC node cannot support a proof edge. An OPEN OBLIGATION blocks descendants that require it. A CLAIM UNDER AUDIT is a document to inspect, not a theorem to import. This monotonicity rule prevents physical pictures or limited calculations from being laundered into an all-data claim through prose.

\section{Certificate Schema for Future Work}
A future analytic certificate for O5 must contain the exact PDE setting, viscosity normalization, solution class, cutoff family, definition of $\mathcal E$, and the range of radii for which the estimate is asserted. It must display every constant in \eqref{eq:decay} or \eqref{eq:depletion} and state whether it depends on initial energy, scale, time, or an unproved critical norm. It must include the pressure decomposition and a bound for each local and far-field piece. A statement that the remainder is negligible without a norm and summation argument is rejected.

If computation enters a future O4 or O7 component, the certificate must identify discretization, precision, stability estimate, data source, and the exact finite domain it covers. It may support a numerical conjecture or finite estimate. It may not assert the continuum all-scale problem unless a theorem transfers the finite certificate to every unresolved scale. A second verifier should consume the certificate rather than reuse the original code path.

For a public breakthrough claim, the audit card records bibliographic identity, version, target formulation, new definitions, imported theorems, hypotheses, every scale argument, treatment of pressure, limiting procedure, and the final universal implication. The card may find a gap without asserting that the claimed conclusion is false. It reports only that the displayed route has not discharged a named obligation.

\section{Conditional Consequence Template}
The intended theorem shape is conditional until O5--O7 are proved. Let $\mathcal E$ be the envelope in \eqref{eq:envelope}. Assume every candidate singular cylinder satisfies \eqref{eq:decay}, $\mathcal T_{\mathrm{rem}}$ is uniformly summable along the geometric scales $\theta^jr$, and constants are independent of the selected cylinder. Then iteration yields
\begin{equation}
\mathcal E(\theta^j r,z_0)
\le q^j\mathcal E(r,z_0)+
K\sum_{\ell=0}^{j-1}q^{j-1-\ell}
\mathcal T_{\mathrm{rem}}(\theta^\ell r,z_0).
\label{eq:iterate}
\end{equation}
If the right side eventually falls below the epsilon-regularity threshold, $z_0$ is regular. If every first singularity is covered, no first singularity exists and the continuation theorem gives global smoothness.

The template makes the missing work visible. This report does not establish the assumptions preceding \eqref{eq:iterate}. It establishes only that they are sufficient in the proposed architecture and that their proof must be pressure-safe, scale-uniform, and independent of the conclusion.

\begin{thebibliography}{00}
\bibitem{fefferman} C. L. Fefferman, "Existence and smoothness of the Navier--Stokes equation," Clay Mathematics Institute, 2000. [Online]. Available: \url{https://www.claymath.org/millennium-problems/navier-stokes-equation}
\bibitem{leray} J. Leray, "Sur le mouvement d'un liquide visqueux emplissant l'espace," \emph{Acta Mathematica}, vol. 63, pp. 193--248, 1934, doi: 10.1007/BF02547354.
\bibitem{ckn} L. Caffarelli, R. Kohn, and L. Nirenberg, "Partial regularity of suitable weak solutions of the Navier--Stokes equations," \emph{Communications on Pure and Applied Mathematics}, vol. 35, no. 6, pp. 771--831, 1982, doi: 10.1002/cpa.3160350604.
\bibitem{ess} L. Escauriaza, G. Seregin, and V. Sverak, "L$_{3,\infty}$-solutions of Navier--Stokes equations and backward uniqueness," \emph{Russian Mathematical Surveys}, vol. 58, no. 2, pp. 211--250, 2003, doi: 10.1070/RM2003v058n02ABEH000608.
\bibitem{kenigkoch} C. E. Kenig and G. S. Koch, "An alternative approach to regularity for the Navier--Stokes equations in critical spaces," \emph{Annales de l'Institut Henri Poincare C}, vol. 28, no. 2, pp. 159--187, 2011, doi: 10.1016/j.anihpc.2010.10.004.
\bibitem{knss} G. S. Koch, N. Nadirashvili, G. Seregin, and V. Sverak, "Liouville theorems for the Navier--Stokes equations and applications," \emph{Acta Mathematica}, vol. 203, no. 1, pp. 83--105, 2009, doi: 10.1007/s11511-009-0039-6.
\bibitem{tao} T. Tao, "Localisation and compactness properties of the Navier--Stokes global regularity problem," arXiv:1108.1165, 2011, doi: 10.48550/arXiv.1108.1165.
\bibitem{barkerprange} T. Barker and C. Prange, "Quantitative regularity for the Navier--Stokes equations via spatial concentration," \emph{Communications in Mathematical Physics}, vol. 385, no. 2, pp. 717--792, 2021, doi: 10.1007/s00220-021-04122-x.
\bibitem{farwig} R. Farwig, "From Jean Leray to the millennium problem: the Navier--Stokes equations," \emph{Journal of Evolution Equations}, vol. 21, no. 3, pp. 3243--3263, 2021, doi: 10.1007/s00028-020-00645-3.
\bibitem{buaria} D. Buaria, A. Pumir, and E. Bodenschatz, "Self-attenuation of extreme events in Navier--Stokes turbulence," \emph{Nature Communications}, vol. 11, p. 5852, 2020, doi: 10.1038/s41467-020-19530-1.
\bibitem{raugel} G. Raugel and G. R. Sell, "Navier--Stokes equations on thin 3D domains. I. Global attractors and global regularity of solutions," \emph{Journal of the American Mathematical Society}, vol. 6, no. 3, pp. 503--568, 1993, doi: 10.1090/S0894-0347-1993-1179539-4.
\bibitem{necas} J. Necas, M. Ruzicka, and V. Sverak, "On Leray's self-similar solutions of the Navier--Stokes equations," \emph{Acta Mathematica}, vol. 176, no. 2, pp. 283--294, 1996, doi: 10.1007/BF02551584.
\bibitem{berselli} L. C. Berselli and G. P. Galdi, "Regularity criteria involving the pressure for the weak solutions to the Navier--Stokes equations," \emph{Proceedings of the American Mathematical Society}, vol. 130, no. 12, pp. 3585--3595, 2002, doi: 10.1090/S0002-9939-02-06697-2.
\end{thebibliography}
\end{document}
